% =====================================================================
%  Bölüm 15 — Bağımlı Tipler ve Lean'in Çekirdeği
% =====================================================================
\chapter{Bağımlı Tipler ve Lean'in Çekirdeği}\label{bol:bagimli}

\begin{bolumhedefleri}
\begin{itemize}
  \item Bağımlı fonksiyon ($Π$) ve bağımlı çift ($Σ$) tiplerini ve bunların
        $\forall$ ve $\exists$ ile ilişkisini,
  \item Tip aileleriyle ``doğruluğu tipine yazılmış'' veri yapıları kurmayı
        (uzunluk indisli vektörler, \lc{Fin}),
  \item Martin-Löf tip kuramını, İnşalar Hesabını ve $\lambda$-küpünü,
  \item \Lean{}'in evren hiyerarşisini ve \lc{Prop}'un impredikatifliğini,
  \item Kanıt-ilgisizliğini, tanımsal ve önermesel eşitlik ayrımını,
  \item \Lean{}'in güvenilir çekirdeğini ve \leankw{\#print} \lc{axioms}
        komutunu.
\end{itemize}
\end{bolumhedefleri}

Bölüm~14'teki basit tiplerde terimler ve tipler iki ayrı dünyada
yaşıyordu: Tipler terimlere bağlı olamıyordu. $A \Rightarrow B$ yazabiliyor,
ama ``$n$ uzunluğunda listeler'' gibi bir tip yazamıyorduk. Bu bölümde iki
dünyayı birleştiriyoruz. Tiplerin terimlere bağlı olabildiği bir dilde,
yani \emph{bağımlı tip kuramında}, Curry--Howard eşlemesi önermeler
mantığından yüklemler mantığına genişler: $\forall$ bir bağımlı fonksiyon,
$\exists$ ise bir bağımlı çift olur. \Lean{}'in kendisi böyle bir dildir.
Bu bölümün sonunda \Lean{}'in çekirdeğinin neye inandığını ve neden ona
güvenebileceğimizi anlatacağız.

% ---------------------------------------------------------------------
\section{$Π$-tipleri: bağımlı fonksiyonlar}\label{sec:pi}

Sıradan bir fonksiyon $f : A \to B$, her $a : A$ için \emph{aynı} $B$
tipinde bir değer döndürür. Bağımlı bir fonksiyonun dönüş tipi ise
argümana göre değişebilir:
\[
  f : \prod_{x : A} B(x)
  \qquad\text{\Lean{}'de}\qquad
  f : (x : A) \to B\;x .
\]
Burada $B : A \to \mathsf{Type}$ bir \emph{tip ailesidir}: Her $x$'e bir
tip atar. Basit bir örnekle başlayalım:

\leandosya[Π-tipleri]{Bolum15}{pi}

\noindent \lc{ornekDeger true} bir doğal sayı, \lc{ornekDeger false} ise
bir metindir. \Cref{fig:pi-sigma-lif}'nın sol tarafı bunu geometrik olarak
gösteriyor. Bir tip ailesini, taban tipin her noktasının üzerinde duran
bir ``lif'' olarak düşünebiliriz. Bağımlı bir fonksiyon, her lif üzerinden
bir nokta seçer.

\begin{figure}[htbp]
\centering
\begin{tikzpicture}[
    lif/.style={draw=#1, fill=#1!8, thick, rounded corners=4pt,
                minimum width=1.6cm, minimum height=2.6cm},
    nokta/.style={circle, fill=#1, inner sep=2.2pt},
  ]
  % --- Π ---
  \node[font=\baslikfont\bfseries, text=anarenk] at (1.6,3.6) {$Π$: her liften bir nokta};
  \node[lif=anarenk] (l1) at (0.6,1.6) {};
  \node[lif=ikincil] (l2) at (2.6,1.6) {};
  \node[font=\baslikfont\scriptsize, text=anarenk] at (0.6,3.05) {\lc{Nat}};
  \node[font=\baslikfont\scriptsize, text=ikincil] at (2.6,3.05) {\lc{String}};
  \node[nokta=teoremframe] (p1) at (0.6,1.9) {};
  \node[nokta=teoremframe] (p2) at (2.6,1.2) {};
  \node[font=\scriptsize, right=1pt of p1] {42};
  \node[font=\scriptsize, right=1pt of p2] {``kırk iki''};
  \draw[very thick, black!60] (-0.2,-0.2) -- (3.4,-0.2);
  \node[nokta=black!70] (b1) at (0.6,-0.2) {};
  \node[nokta=black!70] (b2) at (2.6,-0.2) {};
  \node[font=\scriptsize, below=2pt of b1] {\lc{true}};
  \node[font=\scriptsize, below=2pt of b2] {\lc{false}};
  \node[font=\baslikfont\scriptsize, text=black!55] at (4.0,-0.2) {\lc{Bool}};
  \draw[teoremframe, thick, dashed] (p1) -- (p2);
  % --- Σ ---
  \begin{scope}[xshift=7.2cm]
  \node[font=\baslikfont\bfseries, text=ikincil] at (1.6,3.6) {$Σ$: liflerin ayrık birleşimi};
  \node[lif=anarenk] at (0.6,1.6) {};
  \node[lif=ikincil] at (2.6,1.6) {};
  \node[font=\baslikfont\scriptsize, text=anarenk] at (0.6,3.05) {\lc{Vek Nat 0}};
  \node[font=\baslikfont\scriptsize, text=ikincil] at (2.6,3.05) {\lc{Vek Nat 2}};
  \node[nokta=teoremframe] (q) at (2.6,1.6) {};
  \node[font=\scriptsize, below=2pt of q, align=center] {\lc{⟨2, [7,8]⟩}};
  \draw[very thick, black!60] (-0.2,-0.2) -- (3.4,-0.2);
  \node[nokta=black!70] (c1) at (0.6,-0.2) {};
  \node[nokta=black!70] (c2) at (2.6,-0.2) {};
  \node[font=\scriptsize, below=2pt of c1] {0};
  \node[font=\scriptsize, below=2pt of c2] {2};
  \node[font=\scriptsize, text=black!55] at (1.6,1.6) {$\cdots$};
  \node[font=\baslikfont\scriptsize, text=black!55] at (4.0,-0.2) {\lc{Nat}};
  \end{scope}
\end{tikzpicture}
\caption{Bir tip ailesini, taban tipin her noktası üzerinde duran bir
lif olarak düşünmek. Solda: bir $Π$-tipi elemanı, \emph{her} liften bir
nokta seçer (bir ``kesit''). Sağda: bir $Σ$-tipi elemanı, \emph{bir} lif
ve o liften bir noktadır.}
\label{fig:pi-sigma-lif}
\end{figure}

\noindent Bu bakış açısı Bölüm~6'da $\forall$ için verdiğimiz okumayla
örtüşür: $\forall x,\, P\,x$'in bir kanıtı, her $x$ için $P\,x$'in bir
kanıtını veren bir fonksiyondur. Yani \emph{$\forall$, dönüş tipi
\lc{Prop} olan bir $Π$-tipidir}. \lc{pi} kesitindeki son satır bunu
doğruluyor. \Lean{}'de \lc{∀ n : Nat, P n} ile \lc{(n : Nat) → P n} aynı
terimdir. $A \to B$ ise $B$'nin $x$'e bağlı olmadığı özel durumdur.

% ---------------------------------------------------------------------
\section{Doğruluğu tipine yazılmış veri: vektörler}\label{sec:vektor}

Bağımlı tiplerin programlamadaki en bilinen kullanımı, bir veri yapısının
bir özelliğini tipine yazmaktır. Bölüm~9'daki \lc{Liste} tipinin boş bir
listenin başını almaya çalışınca ne yapacağı belli değildi; \lc{Option}
döndürmek zorundaydık. Uzunluğu tipinde taşıyan bir vektör tipi bu
sorunu ortadan kaldırır:

\leandosya[Uzunluk indisli vektörler]{Bolum15}{vek}

\noindent \lc{Vek α n}, Bölüm~9'da gördüğümüz \emph{tümevarımsal
ailelerin} bir örneğidir: \lc{α} bir \emph{parametre}, \lc{n} ise bir
\emph{indistir}. Kurucular indisi belirler: \lc{bos}'un uzunluğu 0,
\lc{ekle a v}'ninki \lc{v}'ninkinden bir fazladır. Fonksiyonlara
bakalım:

\begin{itemize}
  \item \lc{bas}, yalnızca \lc{Vek α (n + 1)} kabul eder. Boş bir vektörün
        tipi \lc{Vek α 0}'dır ve \lc{0}, \lc{n + 1} biçiminde değildir.
        Bu yüzden \lc{bas}'ın tanımında \lc{bos} durumu yoktur: \Lean{} bu
        durumun \emph{olanaksız} olduğunu kendisi görür.
  \item \lc{fermuar}, iki vektörün aynı uzunlukta olmasını \emph{tipinde}
        ister. Uzunlukları farklı iki vektörü fermuarlamaya çalışmak bir
        çalışma zamanı hatası değil, bir tip hatasıdır.
  \item \lc{birlestir}'in sonuç tipi \lc{Vek α (n + m)}'dir. Tanım,
        \lc{Nat.add}'ın ikinci argümanı üzerinde özyinelemeli olmasından
        yararlanır. \lc{bos} durumunda \lc{n + 0} ile \lc{n}, \lc{ekle}
        durumunda \lc{n + (k + 1)} ile \lc{(n + k) + 1} \emph{tanımsal
        olarak} eşittir.
\end{itemize}

\leandosya[Vektörlerle çalışmak]{Bolum15}{vek_ornek}

\begin{dikkat}[Dikkat: İndis aritmetiği]
\lc{birlestir}'in sonuç tipini \lc{Vek α (m + n)} diye yazsaydık
tanım reddedilirdi. \lc{bos} durumunda \lc{0 + n} ile \lc{n}
\emph{tanımsal} olarak eşit değildir; yalnızca \emph{önermesel} olarak
eşittir (\lc{Nat.zero\_add}). Bağımlı tiplerle programlamanın en sık
karşılaşılan güçlüğü budur. Tip denetleyici yalnızca hesaplayarak görülen
eşitlikleri kabul eder; geri kalanlar için eşitliği açıkça kanıtlayıp
\lc{h ▸ v} ile ``taşımak'' gerekir (\cref{sec:tanimsal}).
\end{dikkat}

\subsection{\lc{Fin}: güvenli indisleme}

\lc{Fin n}, $n$'den küçük doğal sayıların tipidir. Bir sayı ile onun $n$'den
küçük olduğunun kanıtından oluşan bir yapıdır. \lc{Vek α n}'yi \lc{Fin n}
ile indislemek ``dizi sınırı aşıldı'' hatasını tip düzeyinde olanaksız
kılar:

\leandosya[Fin ile indisleme]{Bolum15}{al}

\noindent \lc{al}'ın tanımında \lc{bos} durumunda \lc{Fin 0}'ın elemanı
olmadığı için \lc{Fin.elim0} kullanıyoruz. Özyinelemeli durumda, $i + 1 <
k + 1$ kanıtından $i < k$ kanıtını \lc{Nat.lt\_of\_succ\_lt\_succ} ile
elde ediyoruz. Veri ile kanıt bu fonksiyonda iç içedir.

% ---------------------------------------------------------------------
\section{Bağımlı tipler üzerinde ispat}\label{sec:bagimli-ispat}

Vektörleri ters çeviren bir fonksiyon yazalım ve iki kez ters
çevirmenin vektörü değiştirmediğini ispatlayalım. Bölüm~9'da aynı teoremi
listeler için ispatlamıştık. Burada tipler, ters çevirmenin uzunluğu
koruduğunu zaten garanti ediyor:

\leandosya[Ters çevirme]{Bolum15}{ters}

\noindent İspatlar, listelerdekilerle neredeyse aynıdır. Fark, tümevarımın
artık bir \emph{tümevarımsal aile} üzerinde yapılmasıdır. \leantac{induction}
\lc{v}, \lc{v}'nin uzunluğunu da otomatik olarak genelleştirir. \lc{ekle}
durumunda uzunluk \lc{k + 1}, tümevarım hipotezinde ise \lc{k}'dır.

\leandosya[Eşitlik boyunca taşıma]{Bolum15}{donustur}

\noindent \lc{h ▸ v} ifadesi, \lc{h : n = m} eşitliğini kullanarak \lc{v}'yi
\lc{Vek α n} tipinden \lc{Vek α m} tipine ``taşır''. Bölüm~9'daki \lc{Esit}
tanımının eleme ilkesi tam olarak budur: Eşit şeyler, birbirinin yerine
konabilir.

% ---------------------------------------------------------------------
\section{$Σ$-tipleri: bağımlı çiftler}\label{sec:sigma}

$Σ$-tipleri, $Π$-tiplerinin ikizidir. $\sum_{x : A} B(x)$ tipinin bir
elemanı, bir $a : A$ ile bir $b : B(a)$'dan oluşan bir çifttir. İkinci
bileşenin tipi birinci bileşene bağlıdır. \Lean{}'de üç çeşidi vardır:

\leandosya[Σ-tipleri, alt tipler ve $\exists$]{Bolum15}{sigma}

\begin{table}[htbp]
\centering
\small
\renewcommand{\arraystretch}{1.45}
\rowcolors{2}{tablosatir}{white}
\begin{tabular}{>{\raggedright\arraybackslash}p{3.1cm}
                >{\raggedright\arraybackslash}p{2.3cm}
                >{\raggedright\arraybackslash}p{2.5cm}
                >{\raggedright\arraybackslash}p{3.8cm}}
\rowcolor{tabloust}
\tablobaslik{Yazım} & \tablobaslik{İkinci bileşen} & \tablobaslik{Sonuç} & \tablobaslik{Tanık çıkarılabilir mi?} \\
\lc{(x : α) × β x}     & veri (\lc{Type})  & veri (\lc{Type}) & evet: \lc{.1} \\
\lc{\{x : α // P x\}}  & kanıt (\lc{Prop}) & veri (\lc{Type}) & evet: \lc{.val} \\
\lc{∃ x : α, P x}      & kanıt (\lc{Prop}) & önerme (\lc{Prop}) & hayır (yalnızca bir kanıt içinde) \\
\end{tabular}
\caption{\Lean{}'deki üç $Σ$-tipi. Fark, sonucun hangi evrende yaşadığıdır.}
\label{tab:sigma-cesitleri}
\end{table}

\noindent Tablonun son satırı önemlidir. \lc{∃ n, n > 3}'ün bir kanıtından,
\emph{bir program içinde} tanığı (4'ü) çıkaramayız:

\begin{leanhata}[Lean hata mesajı (özet)]
def tanik (h : ∃ n : Nat, n > 3) : Nat :=
  match h with | ⟨n, _⟩ => n
-- hata: Exists.casesOn yalnızca Prop'a eleme yapabilir
\end{leanhata}

\noindent Bunun nedeni bir sonraki kesimde göreceğimiz
\emph{kanıt-ilgisizliğidir}: Bir önermenin iki kanıtı \Lean{} için ayırt
edilemez. \lc{∃ n, n > 3}'ün biri 4'ü, diğeri 5'i tanık gösteren iki kanıtı
\emph{aynı} sayılır. Kanıttan tanığı çıkaran bir fonksiyon, aynı girdiye
iki farklı cevap vermek zorunda kalırdı. Bu yüzden \lc{Prop}'tan
\lc{Type}'a ``büyük eleme'' genel olarak yasaktır. Bir \emph{ispat içinde}
ise tanığı kullanmak serbesttir (\leantac{obtain} \lc{⟨n, hn⟩ := h}), çünkü
sonuç yine bir önermedir. Bu kısıtlama, Bölüm~4'teki yapıcı/klasik
tartışmasının bir başka yüzüdür: \Lean{}'de varlık kanıtları
\emph{hesaplama içeriği taşımaz}. Hesaplama içeriği isteniyorsa $Σ$ ya da
alt tip kullanılır.

% ---------------------------------------------------------------------
\section{Martin-Löf tip kuramı}\label{sec:mltt}

Bağımlı tiplerin mantık için sistemli bir kuramını İsveçli mantıkçı Per
Martin-Löf 1970'lerde geliştirdi \citep{martinlof1984}. Martin-Löf tip
kuramında (MLTT) her tip oluşturucu dört tür kuralla tanıtılır. Bu kalıp,
Bölüm~5'teki giriş/eleme kurallarının tiplere genişletilmiş hâlidir:

\begin{table}[htbp]
\centering
\small
\renewcommand{\arraystretch}{1.6}
\rowcolors{2}{tablosatir}{white}
\begin{tabularx}{\linewidth}{>{\raggedright\arraybackslash}p{2.4cm} X X}
\rowcolor{tabloust}
\tablobaslik{Kural} & \tablobaslik{$Π$} & \tablobaslik{$Σ$} \\
oluşum   & $A$ tip, $x{:}A \vdash B$ tip $\Rightarrow$ $Π_{x:A} B$ tip & $A$ tip, $x{:}A \vdash B$ tip $\Rightarrow$ $Σ_{x:A} B$ tip \\
giriş    & $x{:}A \vdash b : B$ $\Rightarrow$ $\lambda x.\, b : Π_{x:A} B$ & $a : A$, $b : B[a/x]$ $\Rightarrow$ $\langle a, b\rangle : Σ_{x:A}B$ \\
eleme    & $f : Π_{x:A}B$, $a : A$ $\Rightarrow$ $f\,a : B[a/x]$ & $p : Σ_{x:A}B$ $\Rightarrow$ $p.1 : A$, $p.2 : B[p.1/x]$ \\
hesaplama & $(\lambda x.\, b)\,a \equiv b[a/x]$ \hfill($\beta$) & $\langle a, b\rangle.1 \equiv a$, $\langle a, b\rangle.2 \equiv b$ \\
\end{tabularx}
\caption{Martin-Löf tip kuramında $Π$ ve $Σ$ için kurallar. $\equiv$,
tanımsal eşitliktir.}
\label{tab:mltt}
\end{table}

\noindent MLTT'nin üçüncü temel tipi \emph{özdeşlik tipidir}: $a =_A b$.
Bölüm~9'da \lc{Esit} olarak tanımladığımız tümevarımsal aile budur. Tek
kurucusu $\mathsf{refl} : a = a$'dır, eleme ilkesi ise ``eşitleri
birbirinin yerine koyma''dır (\lc{▸}). Özdeşlik tiplerinin beklenmedik
derecede zengin bir yapısı vardır. Bir $p : a = b$ kanıtını $a$'dan $b$'ye
bir \emph{yol} olarak düşünmek, \emph{homotopi tip kuramına} götürür
\citep{hott2013}.

\begin{tarihce}[Tarihçe: Martin-Löf'ten Lean'e]
Martin-Löf'ün ilk sistemi (1971) $\mathsf{Type} : \mathsf{Type}$ kuralını
içeriyordu. Jean-Yves Girard aynı yıl bu sistemin tutarsız olduğunu
gösterdi \citep{girard1972}. Martin-Löf bunun üzerine evren hiyerarşisini
getirdi. Thierry Coquand ve Gérard Huet 1988'de \emph{İnşalar Hesabını}
(Calculus of Constructions, CoC) tanıttılar \citep{coquand1988}. Christine
Paulin-Mohring ve Coquand bu sisteme tümevarımsal tipleri ekleyerek
\emph{Tümevarımsal İnşalar Hesabını} (CIC) kurdular \citep{paulin1993}.
CIC, Coq'un (bugün Rocq) temelidir. \Lean{}'in tip kuramı da CIC'nin bir
çeşitlemesidir \citep{demoura2021lean4,carneiro2019}.
\end{tarihce}

% ---------------------------------------------------------------------
\section{İnşalar Hesabı ve evrenler}\label{sec:evrenler}

\subsection{$\lambda$-küpü}

Basit tipli $\lambda$-hesabından \Lean{}'e giden yolu, Barendregt'in
$\lambda$-küpü ile görselleştirebiliriz \citep{barendregt1991}. Basit
tiplerde yalnızca \emph{terimlere bağlı terimler} vardır (sıradan
fonksiyonlar). Küpün üç ekseni üç yeni bağımlılık ekler
(\cref{fig:lambda-kup}):

\begin{itemize}
  \item \textbf{Tiplere bağlı terimler} (çok biçimlilik): \lc{id : (α :
        Type) → α → α}. Sistem F \citep{girard1989}.
  \item \textbf{Tiplere bağlı tipler} (tip oluşturucular): \lc{List :
        Type → Type}.
  \item \textbf{Terimlere bağlı tipler} (bağımlı tipler): \lc{Vek α :
        Nat → Type}. $\lambda P$, mantıkta yüklemler.
\end{itemize}

\begin{figure}[htbp]
\centering
\begin{tikzpicture}[
    kose/.style={draw=#1, fill=#1!10, thick, rounded corners=3pt,
                 font=\small, inner sep=4pt, minimum width=1.3cm},
    kenar/.style={-{Stealth[length=2mm]}, thick, black!45},
    x={(3.6cm,0cm)}, y={(0cm,3.0cm)}, z={(1.7cm,1.2cm)},
  ]
  \node[kose=anarenk]     (st)  at (0,0,0) {$\lambda{\to}$};
  \node[kose=anarenk]     (f)   at (0,1,0) {$\lambda 2$ (F)};
  \node[kose=anarenk]     (w)   at (0,0,1) {$\lambda\underline{\omega}$};
  \node[kose=anarenk]     (fw)  at (0,1,1) {$\lambda\omega$};
  \node[kose=ikincil]     (p)   at (1,0,0) {$\lambda P$};
  \node[kose=ikincil]     (p2)  at (1,1,0) {$\lambda P2$};
  \node[kose=ikincil]     (pw)  at (1,0,1) {$\lambda P\underline{\omega}$};
  \node[kose=teoremframe] (coc) at (1,1,1) {\textbf{CoC}};
  \draw[kenar] (st) -- (f);  \draw[kenar] (st) -- (w);  \draw[kenar] (st) -- (p);
  \draw[kenar] (f) -- (fw);  \draw[kenar] (f) -- (p2);
  \draw[kenar] (w) -- (fw);  \draw[kenar] (w) -- (pw);
  \draw[kenar] (p) -- (p2);  \draw[kenar] (p) -- (pw);
  \draw[kenar] (fw) -- (coc); \draw[kenar] (p2) -- (coc); \draw[kenar] (pw) -- (coc);
  \node[font=\baslikfont\scriptsize, text=black!60, anchor=north] at (0.5,-0.12,0) {terimlere bağlı tipler $\to$};
  \node[font=\baslikfont\scriptsize, text=black!60, anchor=south, rotate=90] at (-0.2,0.5,0) {çok biçimlilik $\rightarrow$};
  \node[font=\baslikfont\scriptsize, text=black!60, anchor=north west] at (0.02,-0.02,0.5) {tip oluşturucular};
  \node[font=\baslikfont\scriptsize, text=teoremframe, anchor=west, align=left] at (1.15,1.1,1)
       {+ tümevarımsal tipler\\+ evrenler\\$\Rightarrow$ CIC, \Lean{}};
\end{tikzpicture}
\caption{Barendregt'in $\lambda$-küpü. Sol alt köşe Bölüm~14'ün basit tipli
$\lambda$-hesabıdır. Her eksen bir bağımlılık türü ekler. Sağ üst köşe,
Coquand ve Huet'nin İnşalar Hesabıdır. \Lean{}'in tip kuramı buna
tümevarımsal tipleri ve bir evren hiyerarşisini ekler.}
\label{fig:lambda-kup}
\end{figure}

\subsection{Evrenler}

\lc{Nat}'ın tipi \lc{Type}'tır. Peki \lc{Type}'ın tipi nedir?
\lc{Type : Type} olamaz; Girard paradoksu bunu tutarsız kılar. \Lean{}'in
çözümü sonsuz bir evren hiyerarşisidir (\cref{fig:evrenler}):

\leandosya[Evrenler]{Bolum15}{evrenler}

\begin{figure}[htbp]
\centering
\begin{tikzpicture}[
    evren/.style={draw=#1, fill=#1!8, thick, rounded corners=4pt,
                  minimum width=2.6cm, minimum height=0.8cm, font=\small},
    ok/.style={-{Stealth[length=2mm]}, thick, black!40},
    ornek/.style={font=\scriptsize, text=black!60, anchor=west},
  ]
  \node[evren=teoremframe] (p)  at (0,0)   {\lc{Prop} = \lc{Sort 0}};
  \node[evren=anarenk]     (t0) at (0,1.4) {\lc{Type} = \lc{Sort 1}};
  \node[evren=anarenk]     (t1) at (0,2.8) {\lc{Type 1} = \lc{Sort 2}};
  \node[font=\large, text=black!40] (dd) at (0,3.9) {$\vdots$};
  \draw[ok] (p) -- node[right, font=\scriptsize]{:} (t0);
  \draw[ok] (t0) -- node[right, font=\scriptsize]{:} (t1);
  \draw[ok] (t1) -- (dd);
  \node[ornek] at (1.7,0)   {\lc{2 + 2 = 4}, \lc{∀ p : Prop, p → p}, \dots \quad (önermeler)};
  \node[ornek] at (1.7,1.4) {\lc{Nat}, \lc{List Nat}, \lc{Vek Nat 3}, \lc{Prop}, \dots};
  \node[ornek] at (1.7,2.8) {\lc{Type}, \lc{ChurchSayi}, \lc{(α : Type) → α → α}, \dots};
\end{tikzpicture}
\caption{\Lean{}'in evren hiyerarşisi. Her evren bir üsttekinin elemanıdır.
\lc{Prop} en alttadır ve \emph{impredikatiftir}: \lc{Prop}'lar üzerinde
niceleme yapan bir önerme yine \lc{Prop}'tadır.}
\label{fig:evrenler}
\end{figure}

\noindent İki niceleme arasındaki fark incedir ama önemlidir.
\lc{(α : Type) → α → α} tipi \emph{bütün} tipler üzerinde nicelediği için
\lc{Type}'ın içinde olamaz; bir üst evrene, \lc{Type 1}'e çıkar. Buna
\emph{predikatif} niceleme denir. \lc{∀ p : Prop, p → p} ise \emph{bütün}
önermeler üzerinde niceler, kendisi de bunlardan biri olmasına rağmen
\lc{Prop}'ta kalır. Bu \emph{impredikatif} niceleme, \ciltiki{5}'te göreceğimiz
gibi klasik matematiğin büyük bir kısmını \Lean{}'de ifade etmeyi mümkün
kılar. Tutarlılığı ise \lc{Prop}'tan büyük elemenin kısıtlanmasına bağlıdır
(\cref{sec:sigma}).

\noindent Evren çok biçimliliği (\leankw{universe} \lc{u}) sayesinde
\lc{List} ya da \lc{Vek} gibi tanımlar her evrende bir kez yazılır.
\lc{\#check @List.map}'in çıktısındaki \lc{u\_1}, \lc{u\_2} bu evren
değişkenleridir.

% ---------------------------------------------------------------------
\section{Kanıt-ilgisizliği ve tanımsal eşitlik}\label{sec:tanimsal}

\Lean{}'de \lc{Prop}'un elemanları için \emph{tanımsal kanıt-ilgisizliği}
geçerlidir: Aynı önermenin herhangi iki kanıtı tanımsal olarak eşittir.

\leandosya[Kanıt-ilgisizliği]{Bolum15}{kanit_ilgisiz}

\noindent Bu, Bölüm~9'da \lc{Prop} ile \lc{Bool}'u karşılaştırırken
vurguladığımız ayrımın ta kendisidir. Bir önermede önemli olan
\emph{doğru olup olmadığıdır}, nasıl kanıtlandığı değil. Pratikte bunun
büyük bir yararı vardır: \lc{\{n : Nat // n > 0\}} ya da \lc{Fin n} gibi
tiplerin elemanları yalnızca değerleriyle belirlenir. Kanıt bileşenleri
eşitlik ispatlarında hiç dert olmaz.

\subsection{İki tür eşitlik}

Kitap boyunca iki farklı eşitlik kullandık. Şimdi aralarındaki farkı
kesin olarak söyleyebiliriz:

\begin{table}[htbp]
\centering
\small
\renewcommand{\arraystretch}{1.5}
\rowcolors{2}{tablosatir}{white}
\begin{tabularx}{\linewidth}{>{\raggedright\arraybackslash}p{2.6cm} X X}
\rowcolor{tabloust}
\tablobaslik{} & \tablobaslik{Tanımsal eşitlik ($\equiv$)} & \tablobaslik{Önermesel eşitlik ($=$)} \\
nedir?          & bir \emph{yargı}: iki terim hesaplanınca aynı olur & bir \emph{önerme}: \lc{a = b} bir tiptir \\
kim karar verir? & tip denetleyici, otomatik olarak & biz, bir ispatla \\
örnek           & \lc{n + 0 ≡ n}, \lc{(fun x => x + 1) 2 ≡ 3} & \lc{0 + n = n} (tümevarım gerekir) \\
\Lean{}'de      & \leantac{rfl}, \leantac{show}, \leantac{change}, \leantac{decide} & \leantac{rw}, \leantac{simp}, \lc{▸}, \leantac{calc} \\
karar verilebilir mi? & ilke olarak evet (tip denetimi buna dayanır) & genel olarak hayır \\
\end{tabularx}
\caption{Tanımsal ve önermesel eşitlik.}
\label{tab:esitlikler}
\end{table}

\leandosya[Tanımsal ve önermesel eşitlik]{Bolum15}{tanimsal}

\noindent Tanımsal eşitlik ``bedavadır'': \Lean{} onu görmek için hiçbir
kanıt istemez. Ama onu kullanan bir tanım, tanımın \emph{nasıl} yazıldığına
bağımlı hâle gelir. \lc{n + 0 = n}'nin \leantac{rfl} ile, \lc{0 + n = n}'nin
ise yalnızca tümevarımla ispatlanabilmesinin nedeni, \lc{Nat.add}'ın
ikinci argümanı üzerinde özyinelemeli tanımlanmış olmasıdır.

\begin{dikkat}[Dikkat: Tanımsal eşitliğin sınırları]
Kuramsal olarak \Lean{}'in tanımsal eşitliği karar verilebilir
\emph{değildir}: Kanıt-ilgisizliği ile alt tekil eleme
(\emph{subsingleton elimination}) birleşince, tanımsal eşitliğin geçişli
olmadığı yapay örnekler kurulabilir \citep{carneiro2019}. Uygulamada bu
örneklere neredeyse hiç rastlanmaz. \Lean{}'in çekirdeği tanımsal eşitliği
bir algoritmayla denetler; algoritma bazı durumlarda ``eşit değil''
diyebilir, ama asla yanlış yere ``eşit'' demez. Tutarlılık için önemli
olan da budur.
\end{dikkat}

% ---------------------------------------------------------------------
\section{Lean'in çekirdeği}\label{sec:cekirdek}

Bir ispat asistanına neden güvenelim? Asistanın kendisi binlerce satırlık
bir programdır ve hatalar içerebilir. Nicolaas de~Bruijn'in Automath
projesinden \citep{debruijn1970} kalan cevap, \emph{de~Bruijn ölçütü}
olarak bilinir:

\begin{onemlifikir}[Önemli Fikir: de Bruijn ölçütü]
Bir ispat asistanı, ürettiği her ispatı küçük ve bağımsız olarak
denetlenebilir bir \emph{çekirdeğe} doğrulatmalıdır. Taktikler, otomasyon,
notasyon ve tip sınıfı çözümlemesi ne kadar karmaşık olursa olsun, sonuçta
çekirdeğin anlayacağı basit bir terim üretilir. Güvenmemiz gereken tek
şey, bu terimi denetleyen çekirdektir.
\end{onemlifikir}

\begin{figure}[htbp]
\centering
\begin{tikzpicture}[
    kat/.style={draw=#1, fill=#1!8, thick, rounded corners=5pt,
                text width=3.0cm, align=center, inner sep=6pt, font=\small,
                minimum height=1.6cm},
    ok/.style={-{Stealth[length=2.6mm]}, very thick, black!40},
  ]
  \node[kat=arastirmaframe] (k) at (0,0)   {\textbf{Kaynak kod}\\{\scriptsize \leantac{simp}, \leantac{omega}, notasyon, \lc{x + 1}}};
  \node[kat=anarenk]        (e) at (4.4,0) {\textbf{Elaboratör}\\{\scriptsize taktikler, tip sınıfları, örtük argümanlar}};
  \node[kat=teoremframe]    (c) at (8.8,0) {\textbf{Çekirdek}\\{\scriptsize CIC terimleri, tip denetimi, 3 aksiyom}};
  \draw[ok] (k) -- (e);
  \draw[ok] (e) -- node[above, font=\baslikfont\scriptsize]{terim} (c);
  \node[font=\baslikfont\scriptsize, text=hataframe, anchor=north] at (4.4,-1.05) {güvenilmez (hata yapabilir)};
  \node[font=\baslikfont\scriptsize, text=bittiframe, anchor=north] at (8.8,-1.05) {güvenilir (küçük, denetlenebilir)};
  \node[draw=bittiframe, dashed, thick, rounded corners=4pt, fit=(c), inner sep=8pt] {};
  \node[font=\baslikfont\scriptsize, text=black!55, align=center] at (8.8,1.45) {bağımsız denetleyiciler:\\Lean4Lean, nanoda, \dots};
\end{tikzpicture}
\caption{\Lean{}'in mimarisi. Elaboratörün hataları en kötü ihtimalle
geçerli bir ispatın reddedilmesine yol açar. Geçersiz bir ispatın kabul
edilmesi için çekirdekte bir hata olması gerekir.}
\label{fig:lean-mimari}
\end{figure}

\noindent Elaboratörün ne kadar iş yaptığını görmek için
\lc{pp.all} seçeneğiyle bir terimin çekirdekteki hâline bakabiliriz:

\leandosya[Elaborasyon]{Bolum15}{elaborasyon}

\noindent Masum görünen \lc{x + 1}, aslında üç tip sınıfı örneği ve bir
sayı sabiti kodlaması içerir. Bir taktik ispatı ise çok daha büyük bir
terime dönüşür. Çekirdeğin gördüğü her tümevarımsal tip, kendi
\emph{özyineleyicisi} ile birlikte gelir:

\leandosya[Özyineleyiciler]{Bolum15}{ozyineleyici}

\noindent \lc{Vek.rec}'in tipi Bölüm~9'daki tümevarım ilkelerinin bağımlı
hâlidir. \emph{Motive}, hem indise (\lc{n}) hem de vektörün kendisine
bağlıdır. \leankw{match} ile yazdığımız her tanım ve \leantac{induction}
ile yaptığımız her ispat, çekirdekte bu özyineleyiciye çevrilir.

\subsection{Aksiyomlar}

\Lean{}'in çekirdeği CIC'ye yalnızca üç aksiyom ekler:

\begin{center}
\renewcommand{\arraystretch}{1.45}
\rowcolors{2}{tablosatir}{white}
\begin{tabular}{>{\raggedright\arraybackslash}p{3.4cm}
                >{\raggedright\arraybackslash}p{5.2cm}
                >{\raggedright\arraybackslash}p{3.6cm}}
\rowcolor{tabloust}
\tablobaslik{Aksiyom} & \tablobaslik{Söylediği} & \tablobaslik{Nerede?} \\
\lc{propext}          & denk önermeler eşittir: \lc{(p ↔ q) → p = q} & \leantac{simp}, \leantac{rw} ile ↔ \\
\lc{Quot.sound}       & bölüm tiplerinde ilişkili elemanlar eşittir & fonksiyon dışsallığı (\lc{funext}) \\
\lc{Classical.choice} & boş olmayan her tipten bir eleman seçilebilir & üçüncü hâl, \lc{Classical} \\
\end{tabular}
\end{center}

\noindent Bir teoremin hangi aksiyomlara dayandığını \leankw{\#print}
\lc{axioms} ile öğrenebiliriz:

\leandosya[Aksiyomları denetlemek]{Bolum15}{aksiyomlar}

\noindent Vektör teoremimiz hiçbir aksiyoma dayanmıyor: Tamamen yapıcıdır.
Üçüncü hâlin olanaksızlığı ise üç aksiyomun üçüne de dayanır. Bölüm~4'te
yapıcı ile klasik arasında çizdiğimiz çizgi, \Lean{}'de tam olarak bu
komutla denetlenebilir. Diaconescu teoremi \citep{diaconescu1975},
seçim aksiyomunun \lc{propext} ve \lc{funext} ile birlikte üçüncü hâli
gerektirdiğini söyler. Bu yüzden \lc{Classical.em} bir aksiyom değil, bir
\emph{teoremdir}. Aksiyomları ve onların kuramsal sonuçlarını Bölüm~16'da
ayrıntılı olarak ele alacağız.

% ---------------------------------------------------------------------
\begin{arastirmanotu}
\textbf{Tip kuramının çeşitleri.} \Lean{}, kanıt-ilgisizliği ve
\lc{propext} sayesinde ``klasik matematik dostu'' bir tasarım seçmiştir.
Başka seçimler de mümkündür. \emph{Homotopi tip kuramı} \citep{hott2013},
kanıt-ilgisizliği yerine eşitlik kanıtlarına yollar olarak zengin bir yapı
verir ve \emph{tekdeğerlilik} (\emph{univalence}) aksiyomunu ekler:
Denk tipler eşittir. Kübik tip kuramı bu aksiyomu hesaplanabilir kılar ve
Agda'nın kübik kipinde kullanılır. \emph{Gözlemsel tip kuramı}
(\emph{observational type theory}) ise fonksiyon dışsallığını aksiyom
olmadan tanımsal olarak sağlamayı hedefler.

\textbf{Çekirdeğin doğrulanması.} Mario Carneiro'nun yüksek lisans tezi
\citep{carneiro2019}, \Lean{}'in tip kuramının ilk kesin tanımını verdi ve
tutarlılığını küme kuramında bir modelle gösterdi. Lean4Lean
\citep{carneiro2026}, bu çalışmayı \Lean{}'in içine taşır: \Lean{}'in
çekirdeğinin bir kopyası \Lean{}'de yazılmıştır ve özelliklerinin
ispatları da \Lean{}'dedir. Gödel'in ikinci eksiklik teoremi (\ciltiki{4})
nedeniyle bir sistemin kendi tutarlılığını ispatlayamayacağını
hatırlatalım. Bu yüzden bu tür çalışmalar tutarlılığı ancak daha güçlü bir
sistemde (örneğin ek bir evrenle) ifade edebilir.
\end{arastirmanotu}

\begin{bolumozeti}
\begin{itemize}
  \item $Π$-tipleri bağımlı fonksiyonlardır. $\forall$, dönüş tipi
        \lc{Prop} olan bir $Π$-tipidir. $A \to B$ bağımsız özel
        durumdur.
  \item $Σ$-tipleri bağımlı çiftlerdir. \Lean{}'de üç çeşidi vardır:
        \lc{Sigma} (veri), alt tip (veri + kanıt) ve \lc{Exists} (önerme).
        \lc{Exists}'ten tanık bir program içinde çıkarılamaz.
  \item Tümevarımsal aileler (\lc{Vek α n}, \lc{Fin n}) bir özelliği tipe
        yazarak çalışma zamanı hatalarını tip hatalarına dönüştürür. Bunun
        bedeli indis aritmetiği ve eşitlik boyunca taşımadır.
  \item Martin-Löf tip kuramı her tipi oluşum, giriş, eleme ve hesaplama
        kurallarıyla tanıtır. CIC ve \Lean{}, bu çizginin
        devamıdır. $\lambda$-küpü bağımlılık türlerini düzenler.
  \item Evrenler bir hiyerarşi oluşturur. \lc{Prop} impredikatiftir,
        \lc{Type u}'lar predikatiftir.
  \item \lc{Prop}'ta kanıt-ilgisizliği tanımsaldır. Tanımsal eşitlik tip
        denetleyicinin yargısıdır, önermesel eşitlik ise ispat ister.
  \item de~Bruijn ölçütü: Güven yalnızca küçük çekirdeğe dayanır.
        \Lean{}'in üç aksiyomu vardır ve \leankw{\#print} \lc{axioms} bir
        teoremin hangilerine dayandığını gösterir.
\end{itemize}
\end{bolumozeti}

\begin{alistirmalar}
\alistirma{1} \lc{tekrarlaV a n}'yi \lc{f} ile eşlemek, \lc{tekrarlaV (f a)
n}'yi verir. (\lc{a15\_1})

\alistirma{1} \lc{son (v.sonaEkle a) = a}. \lc{son}'un tanımında neden
uzunluk üzerinde örüntü eşleme yapılıyor? (\lc{a15\_2})

\alistirma{2} Fermuarın birinci bileşenleri ilk vektörü verir. (\lc{a15\_3})

\alistirma{2} \lc{birlestir} ile \lc{toList} uyumludur. (\lc{a15\_4})

\alistirma{2} Eşitlik boyunca taşımak, vektörün listesini değiştirmez.
\leantac{subst}'ün burada ne yaptığını açıklayın. (\lc{a15\_5})

\alistirma{1} \lc{Fin}'de kanıt-ilgisizliği. (\lc{a15\_6})

\alistirma{2} Her liste, uzunluğu kadar bir vektöre dönüşür. Sonuç neden bir
$Σ$-tipi? (\lc{a15\_7})

\alistirma{3} Sabit bir vektörün her elemanı aynıdır. (\lc{a15\_8})

\alistirma{2} (Kâğıt-kalem) \lc{(α : Type) → List α → Nat} ve
\lc{∀ α : Type, Nonempty α → True} tiplerinin hangi evrende yaşadığını
belirleyin ve \leankw{\#check} ile doğrulayın.

\alistirma{2} (Kâğıt-kalem) \lc{birlestir}'in sonuç tipini \lc{Vek α (m +
n)} olarak yazıp \Lean{}'in hata mesajını okuyun. Tanımı
\lc{Nat.add\_comm} ve \lc{▸} ile düzeltin.

\alistirma{3} (Proje) \lc{Vek} için bir \lc{ters} fonksiyonunu
\lc{sonaEkle} kullanmadan, bir biriktirici ile tanımlayın:
\lc{tersAux : Vek α m → Vek α n → Vek α (m + n)}. Hangi indis eşitliğine
ihtiyaç duyuyorsunuz?

\alistirma{3} (Araştırma) Bölüm~14'teki basit tipli $\lambda$-hesabına
bağımlı fonksiyon tipleri ekleyerek $\lambda P$'yi tanımlayın. Tipler artık
terimler içerebildiği için \lc{Tip} ile \lc{Terim} neden aynı tümevarımsal
tip içinde tanımlanmak zorundadır?
\end{alistirmalar}
